% Generated draft: requires mathematical and source review.
\documentclass[11pt]{article}
\usepackage[margin=25mm]{geometry}
\usepackage{fontspec}
\setmainfont{Latin Modern Roman}
\usepackage{amsmath,amssymb,mathtools,unicode-math}
\setmathfont{Latin Modern Math}
\usepackage{enumitem,needspace,xurl,hyperref}
\setlength{\parindent}{0pt}
\setlength{\parskip}{5pt}
\setlength{\emergencystretch}{3em}
\newcommand{\N}{\mathbb{N}}
\newcommand{\Z}{\mathbb{Z}}
\newcommand{\Q}{\mathbb{Q}}
\newcommand{\R}{\mathbb{R}}
\newcommand{\C}{\mathbb{C}}
\newcommand{\F}{\mathbb{F}}
\newcommand{\A}{\mathbb{A}}
\newcommand{\PP}{\mathbb{P}}
\newcommand{\E}{\mathbb{E}}
\newcommand{\eps}{\varepsilon}
\newcommand{\dd}{\,\mathrm{d}}
\begin{document}
\section*{UnsolvedMath ID: 30004757}
\textbf{Source problem code:} catalogue entry
\subsection{Definitions and mathematical statement}
The supplied record concerns singular convex solutions of the displayed equation
\[
\det D^2u=1+\sum_{i\in I}a_i\delta_{p_i},
\]
where \(\delta_{p_i}\) denotes the unit Dirac measure at \(p_i\). The dimension, domain, index set, coefficients, boundary conditions, and particular family of examples are not specified.

A standard conditional interpretation is the Aleksandrov Monge--Amp\`ere equation. To describe this interpretation, let \(\Omega\subset\R^n\) be open and convex and let \(u:\Omega\to\R\) be finite and convex. Define
\[
\partial u(x)=\{q\in\R^n:u(y)\ge u(x)+q\cdot(y-x)\text{ for every }y\in\Omega\},
\]
and, for Borel sets \(E\subset\Omega\), define its Monge--Amp\`ere measure by
\[
\operatorname{MA}_u(E)=\mathcal L^n\!\left(\bigcup_{x\in E}\partial u(x)\right),
\]
where \(\mathcal L^n\) is Lebesgue measure. In this interpretation the displayed equation means
\[
\operatorname{MA}_u=\mathcal L^n\!\restriction_\Omega+\sum_{i\in I}a_i\delta_{p_i}
\quad\text{as measures on }\Omega,
\]
provided the atomic expression defines a nonnegative locally finite measure. This replaces the scalar \(1\) by Lebesgue measure and interprets the determinant in a measure sense. These are explicit interpretive conventions, not hypotheses or a solution concept confirmed by the supplied record.

For an interior point \(p\) of such a convex function, a standard first-order tangent profile is
\[
T_pu(v)=\lim_{t\downarrow0}\frac{u(p+tv)-u(p)}{t},\qquad v\in\R^n.
\]
Its graph \(\{(v,T_pu(v)):v\in\R^n\}\) describes a graph tangent cone translated to the origin; the corresponding epigraph cone is \(\{(v,s):s\ge T_pu(v)\}\). The source does not specify which convention, or which normalization of tangent cones, is intended.

\textbf{Problem, with unresolved scope.} For the source-intended singular convex Monge--Amp\`ere examples satisfying the displayed equation, determine precise descriptions of the tangent cones at each point \(p_i\) carrying a Dirac mass. The original statement additionally asks for analogous investigations when the obstacle graphs are not necessarily convex polytopes. The obstacles, their relation to \(u\), and the admissible replacement class are not supplied, so this additional target cannot be formulated more precisely from the record. It must be retained even though it is absent from the clean statement. No universal quantification over all convex solutions, prescribed dimension, boundary value problem, or obstacle construction is justified by the supplied text.
\subsection{Short English statement}
Describe the tangent cones at the atomic singularities of the intended convex Monge--Amp\`ere examples, and investigate the analogous question for obstacle graphs that need not be convex polytopes. The examples, solution concept, cone convention, and obstacle hypotheses require clarification. The supplied open status is dated curation, not an independently verified conclusion.
\subsection{Sources}
\begin{itemize}
\item Supplied local UnsolvedMath registry record, permanent ID \(30004757\), title: \emph{Dirac-Mass Tangent Cones in Monge--Amp\`ere Equations}; exact source problem code: \texttt{catalogue entry}. Collection: \emph{Oberwolfach Reports Open Problems}; category: Partial Differential Equations. The record gives proposed year \(2021\) and supplies no individual proposer.
\item Supplied source citation: \emph{Partial Differential Equations (2022)}, Oberwolfach Reports, DOI \texttt{10.4171/owr/2021/35}, pp.~1859--1914. \url{https://doi.org/10.4171/owr/2021/35}. The background describes these reports as records of mathematical workshops and their open-problem discussions. The citation's year \(2022\), the DOI's \(2021\) component, and the proposed year are retained as supplied.
\item Supplied literature source: Mooney--Rakshit, \url{https://arxiv.org/abs/2204.11365}. No further bibliographic details or theorem statements are supplied.
\item Supplied literature source: Oberwolfach report, \url{https://ems.press/content/serial-article-files/46911}.
\item Literature assessment dated exactly \texttt{2026-08-22}: the record labels the problem open and states that the higher-dimensional constructions do not give a precise classification of tangent cones at the Dirac masses in the requested examples. This is supplied dated curation derived from the OWR material, not independently verified here; the record explicitly requests independent verification of time-sensitive status and cited claims.
\item The supplied statement status is \texttt{corrected\_verified}. Its verification field reports that the official report was checked directly and that the clean statement preserves the complete mathematical target while restoring definitions and removing extraction artifacts or source-dependent wording. This is a supplied verification claim, not a source visit performed here. The original statement's additional obstacle question is absent from the clean statement and remains unresolved.
\end{itemize}
\end{document}
